\section{実装}

\subsection{実装基盤と全体構成}
実装はGhidraおよびPyGhidraを用いてHigh P-codeを取得する．システムは，解析部，等価性判定部，SMT式生成部，判定不能理由分類部から構成される．
% TODO(draft): 実コードに対応するモジュール構成図と処理フローを追加する．

\subsection{振る舞い要約とSMT検証}
RETURNまでのデータフローを追跡し，戻り値を表す式を構築する．分岐条件，PHI，出力ポインタへの書き込みを要約し，対象の2つのバイナリから得た式をSMTで検証する．また，コンパイラおよびアーキテクチャの差を扱うため，除算・剰余の正規化，型推論の疑義に対する修正リトライ，ABI差分に対するリトライを行う．
% TODO(draft): 各正規化・リトライの発動条件と健全性を実コードおよびテストに基づいて記述する．

\subsection{判定不能となるケースの扱い}
CALL，CALLIND，メモリ効果，ループは，解析可能な範囲でのみ要約する．一般的なループ要約は現時点で対象外である．解析上の証明責務を解消できない場合はUNKNOWN，ソルバーが打ち切りとなった場合はTIMEOUTとして記録する．

判定不能理由は，coverage不足，未解決の呼び出し・副作用，メモリ同一性不明，ABIまたは型の疑義，ループ要約不足，SMTタイムアウトなどに分類する．
% TODO(draft): 実装が出力する正規のreason IDと，本節の分類との対応表を追加する．
